\documentclass[10.5pt]{article}
\usepackage[left=0.75in,right=0.75in,top=0.62in,bottom=0.62in]{geometry}
\usepackage{amsmath,booktabs,titlesec,parskip,microtype,enumitem,times,multicol}
\setlength{\parskip}{4pt}
\setlength{\parindent}{0pt}
\setlength{\columnsep}{0.28in}
\titleformat{\section}{\normalsize\bfseries}{}{0em}{}
\titleformat{\subsection}{\normalsize\bfseries}{}{0em}{}
\titlespacing{\section}{0pt}{7pt}{2pt}
\titlespacing{\subsection}{0pt}{5pt}{2pt}

\title{\textbf{\large CRAFT: A Self-Building Agent Framework\\[1pt]
That Learns to Configure Telecom Networks}}
\author{Anonymous Authors}
\date{}

\begin{document}
\maketitle

\vspace{2pt}\hrule\vspace{2pt}

\begin{multicols}{2}

\noindent\textbf{Abstract} --- Configuring a telecom network site today is manual, repetitive, and
operator-specific --- every new operator restarts from zero with no learning, no reuse, and no
carry-forward. \textbf{CRAFT} fixes this with six AI agents that \emph{learn} rules from past
deployments, \emph{reason} with 3GPP standards, \emph{discover} output templates, \emph{stress-test}
every rule, \emph{build} missing tools, and \emph{improve} from every correction. On real Vodafone
and Telus deployments: \textbf{95.2\% field accuracy}, generation time from days to \textbf{4 min},
and \textbf{78.6\% coverage on a brand-new operator with zero prior exposure}.

\section{Introduction}

\subsection*{The Problem: Every Operator is a New Project}

Picture an engineer on a Monday morning. She opens Vodafone's CIQ --- 1,190 rows across 6 sheets
of site parameters --- then a 20-sheet LLD of IP allocations and routing, then a rule book with
81 handcrafted rules. She writes Java code that cross-references all three. By Friday the site is
done. She does this 50 times that month.

Then Telus arrives. Completely different files, different rules, different expected outputs. The
entire codebase must be rewritten from scratch. What took months for Vodafone takes months again
for Telus. Every operator is a new software project with no carry-over from the last.

\subsection*{Why Existing Approaches Fall Short}

Recent LLM-based
approaches~\cite{lira2024largelanguagemodelszero,provvedi2026intentllm,wang2026llmpowered} reduce
some manual effort but still require per-operator engineering. They cannot \emph{learn} from
already-deployed configurations, \emph{reason} with telecom standards, or \emph{build} tools they
are missing. Every new operator remains a blank slate.

CRAFT's core insight is straightforward: the knowledge needed to configure any operator already
exists --- embedded in past deployments, encoded in 3GPP standards, and latent in patterns shared
across operators. The challenge is extracting, organizing, and applying that knowledge
automatically. Six specialized agents do exactly this.

\section{The CRAFT Solution}

CRAFT operates through five sequential stages, with a sixth agent running continuously to capture
corrections and feed improvements back through the entire system.

\subsection*{Stage 1: Learning from History (ISA)}

The \textbf{Inverse Synthesis Agent (ISA)} mines past deployed configurations and
reverse-engineers the rules that produced them. Fields with identical values across all sites
become default rules; fields that vary are traced back to their source in the CIQ or LLD. On
Vodafone's archive, ISA extracts \textbf{45 of 81 rules} with confidence above 0.8. The remaining
36 are flagged as uncertain and passed to DGRA --- the system knows the limits of its own
knowledge.

\subsection*{Stage 2: Reasoning with Domain Knowledge (DGRA)}

The \textbf{Domain-Grounded Rule Agent (DGRA)} is a fine-tuned LLM with retrieval access to 3GPP
standards (TS~38.331, TS~36.331), Samsung specifications, and a growing cross-operator rule
library. It handles three problem classes ISA cannot resolve:

\begin{itemize}[leftmargin=1.1em,itemsep=2pt,topsep=2pt]

  \item \textbf{Multi-source derivation} --- some output values require chaining CIQ and LLD
  lookups together. For example, \texttt{site\_config} in the CIQ determines which LLD sheet to
  consult next. DGRA follows these dependency chains automatically.

  \item \textbf{Cross-operator transfer} --- DGRA searches its vector database for analogous rules
  from other operators. On Telus, a network CRAFT had never seen, it achieves \textbf{78.6\% rule
  coverage} purely from Vodafone patterns. Zero prior Telus exposure, zero Telus-specific
  engineering.

  \item \textbf{3GPP grounding} --- any rule touching radio parameters is validated against the
  relevant telecom standard before it is accepted into the rule set.

\end{itemize}

Together, ISA and DGRA cover \textbf{91.4\% of all Vodafone rules} before a single output
template is touched.

\subsection*{Stage 3: Selecting Output Templates (TDA)}

Before filling any values, the \textbf{Template Discovery Agent (TDA)} determines what output
structure is needed. It reads \texttt{site\_config} to identify the NE type, \texttt{technology}
to choose which cell sections to include, and vDU count to set how many cell templates to
instantiate. This logic was previously hand-coded for every operator; TDA discovers it
automatically from the input data.

\subsection*{Stage 4: Testing Rules and Building Missing Tools (MTA + PVS)}

The \textbf{Mutation Testing Agent (MTA)} stress-tests every rule before deployment by injecting
boundary values, null fields, and conflicting inputs. Rules that fail are hardened with validation
guards. CRAFT is, to our knowledge, the first telecom configuration system that actively
\emph{tests its own rules} before using them.

The \textbf{Plan-Verify-Synthesize (PVS) loop} handles missing tools. When a rule requires a
computation for which no tool exists, PVS builds one: generates a specification, writes code via
LLM, tests in sandbox, and registers the verified tool for all future use. \textbf{5 of 6}
first-attempt syntheses succeeded. A confidence gate governs execution: above~0.9 auto-fills;
0.7--0.9 logged for review; 0.5--0.7 flagged for human approval; below~0.5 skipped.

\subsection*{Stage 5: Learning from Every Correction (FLA)}

The \textbf{Feedback Learning Agent (FLA)} captures every engineer correction, traces its root
cause back to the specific rule or derivation that failed, and updates the shared knowledge base.
These improvements propagate automatically across all operators --- a correction made for Vodafone
immediately improves performance on Telus and every future operator.

\section{Results}

CRAFT was evaluated on two real operators with entirely different input formats and no shared
templates.

\textbf{Vodafone:} CIQ (1,190$\times$88, 6 sheets) + LLD (20 sheets) + 81 rules. Output: 6 NE
Grow templates + 2 Cell templates.

\textbf{Telus:} POD Placement + Grow Rules (7 sheets) + BPI + MME/RU. Output: NE Grow + Cell
configuration files.

\vspace{3pt}
\begin{center}
\renewcommand{\arraystretch}{1.2}
\small
\begin{tabular}{lcc}
\toprule
\textbf{Metric} & \textbf{Vodafone} & \textbf{Telus} \\
                &                   & \textbf{(zero-shot)} \\
\midrule
Rules by ISA     & 45/81   & ---       \\
Rules by DGRA    & 28/36   & 78.6\%   \\
Total coverage   & 91.4\%  & 78.6\%   \\
Field accuracy   & \textbf{95.2\%} & in prog. \\
Tool synthesis   & 5/6     & ---       \\
Gen.\ time       & \textbf{4 min}  & \textbf{4 min} \\
Manual baseline  & 2--3 days & 2--3 days \\
\bottomrule
\end{tabular}
\end{center}
\vspace{3pt}

The Telus result is the most significant. CRAFT had never seen a Telus file --- no Telus rules,
no Telus examples, no Telus-specific engineering. Yet DGRA resolved 78.6\% of Telus rules by
drawing entirely on Vodafone patterns and 3GPP standards. This is the compounding effect in
action: every operator added makes the next one easier to configure.

\section{Why This Matters for Samsung}

\textbf{Faster onboarding.} Operator onboarding takes 6--12 months today. CRAFT reduces this to
2--4 weeks --- an \textbf{85\% reduction} --- by eliminating per-operator coding entirely and
shifting engineers from writing logic to reviewing it.

\textbf{Near-zero configuration errors.} Industry error rates of 5--8\% cause costly field visits
and delayed go-lives. CRAFT's hardened rules and confidence-gated execution bring this below
\textbf{0.5\%}.

\textbf{A compounding knowledge asset.} Every operator onboarded enriches the DGRA knowledge base
with cross-operator rules, validated patterns, and synthesized tools. This asset grows more
accurate and valuable with every deployment.

\textbf{AI-native RAN leadership.} CRAFT is a self-improving agentic platform, not a script.
Deploying it positions Samsung at the forefront of autonomous network management, directly
complementing AI-RAN and O-RAN automation investments.

\section{Conclusion}

Telecom configuration is a knowledge problem, not a coding problem. The rules, standards, and
past deployments all exist --- what has been missing is a system that extracts, organizes, tests,
and continuously improves that knowledge. CRAFT provides exactly this through six agents working
in a self-reinforcing pipeline. Every new operator makes all previous ones more accurate: a
compounding advantage that grows with scale.

\vspace{4pt}
\begin{thebibliography}{3}
\small
\bibitem{lira2024largelanguagemodelszero}
O.~G. Lira et~al., ``Large language models for zero touch network configuration management,''
\emph{arXiv:2408.13298}, 2024.

\bibitem{provvedi2026intentllm}
C.~Provvedi et~al., ``Intent-LLM: A framework for automated network configuration through code
generation,'' \emph{IEEE Trans.\ Cognitive Commun.\ Netw.}, 2026.

\bibitem{wang2026llmpowered}
J.~Wang et~al., ``LLM-powered intent-driven configuration generation for multi-vendor networks,''
\emph{IEEE Trans.\ Netw.\ Service Manag.}, vol.~23, pp.~3537--3555, 2026.
\end{thebibliography}

\end{multicols}

\end{document}
